\chapter{Methods}\label{methods}
TODO % intro to this section

% dataset construction for retr + gen evaluation
\section{Dataset}
TODO \cite{He2024HDT}

% ScholarCopilot - architecture only!
\section{ScholarCopilot}\label{scholarcopilot}
Instead of treating retrieval and generation as separate processing stages (see \refsec{introduction}) ScholarCopilot \cite{wang2025scholarcopilottraininglargelanguage} integrates the information retrieval directly into the generation process. The model is trained to generate \emph{retrieval tokens} (\mono{[RET]}), that upon being generated, cause the model to interupt the generation process to retrieve a reference from the retrieval corpus \cite{wang2025scholarcopilottraininglargelanguage}. The ScholarCopilot retrieval corpus\footnote{\url{https://huggingface.co/datasets/TIGER-Lab/ScholarCopilot-Data-v1}} consists of abstracts of scientific papers, which are encoded by the ScholarCopilot model prior to training/inference and stored as a HNSW \cite{malkov2016efficientrobustapproximatenearest} index using Faiss \cite{douze2024faiss} for efficient retrieval (see \refsec{hnsw}). Standard ScholarCopilot uses the retrieved abstract to inform its subsequent generation by replacing the \mono{[RET]} tokens in the generated context with the corresponding abstracts. 

ScholarCopilot is intended to assist with academic writing by providing "[...] text completion and citation suggestions" (see the official demo\footnote{\url{https://huggingface.co/spaces/TIGER-Lab/ScholarCopilot}}\textsuperscript{,}\footnote{\url{https://www.youtube.com/watch?v=QlY7S52sWDA}} for an example). In the following we will thus assume that the model is used to generate a scientific paper with citations. Let
\begin{align}\label{eq:initial-context}
    Y^\mono{in}=y^\mono{in}_1,...,y^\mono{in}_l
\end{align}
denote the initial context given to the model. Since we seek to generate a scientific document, this context may consist of the title and abstract of the paper the model is supposed to generate. Let furthermore 
\begin{align}\label{eq:generated-tokens}
    Y^\mono{gen}_{1:i-1}=y^\mono{gen}_{1},...,y^\mono{gen}_{j-1},c_j,y^\mono{gen}_{j+1},...,y^\mono{gen}_{i-1}
\end{align}
denote a token sequence generated by the model where $c_j$ signifies that a \mono{[RET]} token was generated by the model at generation step $j$. Note that the model may generate the retrieval token at multiple steps, and that the number of generated retrieval tokens is equal to the number of citations in the generated paper. The representation $\hat{c}_j$ of a \mono{[RET]} token $c_j$ computed by the model serves as an aggregated representation of the token sequence generated up to the \mono{[RET]} token (e.g. $Y^\mono{gen}_{1:j-1}$ in \refeq{eq:generated-tokens}), in the same manner as BERT \cite{devlin2019bertpretrainingdeepbidirectional} (and also for example CF-ViT \cite{chen2022cfvitgeneralcoarsetofinemethod}) uses the \mono{[cls]} token as a representation of its input sequence. Thus $\hat{c}_j$ is unique for every $j$ since it is dependent of the corresponding left-hand context. 

When the ScholarCopilot generator $p_G$ generates a retrieval token (that is, $p_G(\ast|Y^\mono{in}Y^\mono{gen}_{1:i-1})$ is maximal for the \mono{[RET]} token, thus $\ast\leftarrow c_i$), the generation process is interrupted and instead, the representation $\hat{c}_i$ is passed to the retriever $p_R(z_i|\hat{c}_i)$ used to retrieve an encoded reference $z_i$ corresponding to an abstract
\begin{align}
    Y^\mono{ref}_i=y^\mono{ref}_{i,1},...,y^\mono{ref}_{i,k}
\end{align}
from the retrieval corpus. Before passing the newly generated context to the generator the retrieval token $c_i$ is replaced with the corresponding abstract to obtain an augmented generation output
\begin{align}
    \mathbb{Y}^\mono{gen}_{1:i}=y^\mono{gen}_{1},...,y^\mono{gen}_{i-1},y^\mono{ref}_{i,1},...,y^\mono{ref}_{i,k}
\end{align}
that is used to inform the sebsequent generation steps with the retrieved information. The generator $p_G(y_{i+1}|Y^\mono{in}\mathbb{Y}^\mono{gen}_{1:i})$ predicts the next token $y_{i+1}$ at generation step $i+1$ based on the given initial context $Y^\mono{in}$ and the augmented output $\mathbb{Y}^\mono{gen}_{1:i}$. See \refalg{alg:scholarcopilot} for a detailed overview of one generation step.

\begin{algorithm}[t]
    \captionsetup{labelfont={sc,bf}, labelsep=newline}
    \KwIn{The initial context $Y^\mono{in}$ and the previously generated sequence $Y^\mono{gen}_{1:i-1}$}
    \KwOut{The initial context $Y^\mono{in}$ and the new generated sequence $Y^\mono{gen}_{1:i}$}
    \TitleOfAlgo{ScholarCopilotGenerationStep($Y^\mono{in},Y^\mono{gen}_{1:i-1}$)}
    $y_i\leftarrow\argmax\limits_{y_i\in Y}p_G(y_i|Y^\mono{in}Y^\mono{gen}_{1:i-1})$\\
    \If{$y_i=\mono{[RET]}$}{
        $\hat{c}_i\leftarrow\mathcal{G}(y_i)$\\
        $z_i\leftarrow\argmax\limits_{z_i\in Z}p_R(z_i|\hat{c}_i)$\\
        $Y^\mono{ref}_i\leftarrow$ getAbstract$(z_i)$\\
        $\mathbb{Y}^\mono{gen}_{1:i}\leftarrow Y^\mono{gen}_{1:i-1}Y^\mono{ref}_i$\\
        \Return{$Y^\mono{in}$, $\mathbb{Y}^\mono{gen}_{1:i}$}
    }
    \Else{
        $Y^\mono{gen}_{1:i}\leftarrow Y^\mono{gen}_{1:i-1}y_i$\\
        \Return{$Y^\mono{in}$, $Y^\mono{gen}_{1:i}$}
    }
    \caption{Overview of one generation step of ScholarCopilot. Here $Y$ denotes the vocabulary (the set of all tokens), $Z$ the retrieval corpus (all encoded abstracts), $\mathcal{G}$ the ScholarCopilot backbone model (see \refsec{qwen}) and getAbstract$()$ a function that returns the unencoded abstract corresponding to $z_i$}\label{alg:scholarcopilot}
\end{algorithm}



% FAISS + HNSW
\subsection{HNSW}\label{hnsw}
TODO

\subsection{ScholarCopilot Training}\label{sc-training}
TODO

% Qwen 2.5 (as the reranker)
\section{Passage Reranking}\label{passage-reranking}
The goal of the \emph{Passage Retrieval} module (PR) introduced with this thesis is it to enhance the citation accuracy and generation quality of ScholarCopilot. As outlined in \refsec{scholarcopilot}, ScholarCopilot represents 

% Batched Self-Consistency
\subsection{Batched Self-Consistency}\label{batched-self-consistency}
TODO

% Min-Max Normalization of scores
\subsection{Min-Max Normalization}\label{normalization}
TODO

